\begin{filecontents}[overwrite]{\jobname.bib}
@phdthesis{wheeler2010ddc,
  author = {David A. Wheeler},
  title = {Fully Countering Trusting Trust through Diverse Double-Compiling (DDC)},
  school = {George Mason University},
  year = {2010},
  url = {https://dwheeler.com/trustingtrust/}
}
@online{mccune2009prover9,
  author = {William McCune},
  title = {Prover9 and Mace4},
  year = {2009},
  organization = {University of New Mexico},
  url = {https://www.cs.unm.edu/~mccune/prover9/}
}
@article{thompson1984trust,
  author = {Ken Thompson},
  title = {Reflections on Trusting Trust},
  journal = {Communications of the ACM},
  volume = {27},
  number = {8},
  pages = {761--763},
  month = aug,
  year = {1984},
  publisher = {ACM},
  doi = {10.1145/358198.358210},
  url = {https://doi.org/10.1145/358198.358210}
}
@techreport{karger1974multics,
  author = {Paul A. Karger and Roger R. Schell},
  title = {Multics Security Evaluation: Vulnerability Analysis},
  institution = {Air Force Electronic Systems Division},
  number = {ESD-TR-74-193},
  address = {Hanscom AFB, Bedford, MA, USA},
  year = {1974},
  url = {https://multicians.org/project-mac/ESD-TR-74-193.pdf}
}
@inproceedings{karger2002thirty,
  author = {Paul A. Karger and Roger R. Schell},
  title = {Thirty Years Later: Lessons from the Multics Security Evaluation},
  booktitle = {Proceedings of the 18th Annual Computer Security Applications Conference (ACSAC)},
  pages = {119--126},
  year = {2002},
  publisher = {IEEE Computer Society},
  doi = {10.1109/CSAC.2002.1176285},
  url = {https://doi.org/10.1109/CSAC.2002.1176285}
}
@book{kunen1980set,
  author    = {Kenneth Kunen},
  title     = {Set Theory: An Introduction to Independence Proofs},
  series    = {Studies in Logic and the Foundations of Mathematics},
  volume    = {102},
  publisher = {North-Holland},
  year      = {1980}
}
@article{heyting1930regeln,
  author  = {Arend Heyting},
  title   = {Die formalen {R}egeln der intuitionistischen {L}ogik},
  journal = {Sitzungsberichte der Preussischen Akademie der Wissenschaften, Physikalisch-mathematische Klasse},
  pages   = {42--56},
  year    = {1930}
}
@book{hesseling2003gnomes,
  author    = {Dennis E. Hesseling},
  title     = {Gnomes in the Fog: The Reception of {B}rouwer's Intuitionism in the 1920s},
  publisher = {Birkh{\"a}user},
  year      = {2003},
  doi       = {10.1007/978-3-0348-8007-7}
}
@online{bootstrappableorg,
  author    = {{Bootstrappable Builds Project}},
  title     = {Bootstrappable Builds},
  year      = {2026},
  url       = {https://bootstrappable.org/}
}
\end{filecontents}

\documentclass[sigconf]{acmart}

\setcopyright{none}
\acmYear{2026}
\acmMonth{06}

\title{Assumption Independence in Diverse Double-Compiling (DDC)}

\author{David A. Wheeler}
\orcid{0000-0003-3115-4122}
\affiliation{%
  \department{Open Source Security Foundation (OpenSSF)}
  \institution{The Linux Foundation}
  \city{San Francisco}
  \state{CA}
  \country{USA}
}
\email{dwheeler@linuxfoundation.org}

\begin{abstract}
Diverse Double-Compiling (DDC) counters the classic ``trusting trust'' attack by using a second, trusted compiler to verify the compiler under test. DDC has been formally proved, and its assumptions have been proven consistent. However, the DDC assumptions have not been shown to be independent. This is important to have confidence that the efforts to verify each DDC assumption are necessary. This paper shows that the DDC proof assumptions are independent.
\end{abstract}

\keywords{Diverse Double-Compiling, DDC, Trusting Trust, Supply Chain Security, Open Source Software, formal proof, assumption independence}

\begin{CCSXML}
<ccs2012>
   <concept>
       <concept_id>10002978.10003022.10003023</concept_id>
       <concept_desc>Security and privacy~Software and application security</concept_desc>
       <concept_significance>500</concept_significance>
       </concept>
   <concept>
       <concept_id>10011007.10011074.10011099.10011102.10011103</concept_id>
       <concept_desc>Software and its engineering~Software verification and validation</concept_desc>
       <concept_significance>500</concept_significance>
       </concept>
   <concept>
       <concept_id>10011007.10011006.10011041</concept_id>
       <concept_desc>Software and its engineering~Compilers</concept_desc>
       <concept_significance>300</concept_significance>
       </concept>
 </ccs2012>
\end{CCSXML}

\ccsdesc[500]{Security and privacy~Software and application security}
\ccsdesc[500]{Software and its engineering~Software verification and validation}
\ccsdesc[300]{Software and its engineering~Compilers}

\begin{document}

\maketitle

\section{Introduction}

Many software programs have unintentional vulnerabilities, but in principle these can be found and fixed using various approaches including human review, traditional static analysis, traditional dynamic analysis, and artificial intelligence (AI). Intentional vulnerabilities in source code can often be found by many of the same techniques. An executable of an ordinary program (as opposed to a program for processing programs) might not correspond to its source code, for example, because its build was subverted or it was subverted after the build. Processes that protect the build can reduce the risk of the first. In addition, if the software has a reproducible build, these subversions can be reliably detected by rebuilding the executable and determining if the resulting executable is the same. All of these measures fail to counter an old and well-known attack.

An Air Force evaluation of Multics by Karger and Schell \cite{karger1974multics} \cite{karger2002thirty}, and Ken Thompson’s Turing award lecture (“Reflections on Trusting Trust”) \cite{thompson1984trust}, showed that compilers can be subverted to insert malicious Trojan horses into critical software, including themselves, as part of the compilation process.  If this “trusting trust” attack goes undetected, even complete analysis of a system’s source code will not find the malicious code that is running.

For decades, the industry treated this as an unresolvable problem. If you can't trust the compiler executable used to build your tools, you can't trust anything built with it, including other compiler executables. This is especially important for open source software, where people verify source code but in practice people execute compiled binaries.

Diverse Double-Compiling (DDC) counters this attack by using a second, independent compiler to compile the parent of the compiler under test, then using that result to compile the source code of the compiler under test \cite{wheeler2010ddc}. If the output of this process matches the original compiler executable, then the executable corresponds to its source code. DDC does not require that arbitrary different compilers produce the same executable output given the same input.  Indeed, this would be extremely unlikely in the real world. DDC instead requires determinism by the parent compiler of the compiler under test. It is generally desirable that a compiler be deterministic, and many real-world compilers are deterministic or can be made deterministic via a flag.

Another countermeasure against the trusting trust, that has more recently received attention, are bootstrappable builds.
\href{http://bootstrappable.org/}{\path{http://bootstrappable.org/}}
Bootstrappable builds work to minimize the amount of bootstrap binaries which are then used to build everything else.
This is a promising approach, but challenging, and many environments
are not based on bootstrappable builds.
DDC provides a great complement to bootstrappable builds.
bootstrappable builds tries to limit opaque binaries, and
DDC can use those bootstrappable builds to verify other binaries.

DDC has been formally (mathematically) proven to produce this result when a set of reasonable assumptions is true \cite{wheeler2010ddc}. The assumptions and goals were defined in First Order Logic (FOL) and generating proofs with prover9 \cite{mccune2009prover9}. There is strong confidence that these proofs are correct; they were automatically generated, checked step-by-step by a different formally-proven tool, and reviewed by multiple people. The same source also proved that these assumptions were consistent (that is, they did not contradict each other).

In actuality, there are three proofs involving DDC:

\begin{enumerate}
	\item Goal \path{source_corresponds_to_executable} (5 assumptions): This is the key proof for DDC.  It shows that given certain assumptions, if stage2 (the result of the DDC process) and cA (the original compiler-under-test) are equal, then the executable cA and the source code sA exactly correspond.
	\item Goal \path{always_equal} (9 assumptions): This proves that, under “normal conditions” (such as when compiler executables have not been rigged and thus do correspond to their respective source code), cA and stage2 are in fact always equal.  Thus, the first proof is actually useful, because its assumptions will often hold.  This also implies that if cA and stage2 are not equal, then at least one of its assumptions is not true.
	\item Goal \path{cP_corresponds_to_sP} (3 assumptions): The previous “\path{always_equal}” proof does not require that a “grandparent” compiler exist, but having one is a common circumstance.  This third proof shows that if there is a grandparent compiler, one of the assumptions of proof \#2 can be proved given other assumptions that may be easier to verify (potentially making DDC even easier to apply in this common case).
\end{enumerate}

\section{Axiomatic Independence in First-Order Logic}

If a proof uses a set of assumptions, that does \emph{not} mean
that the assumptions are independent.
It is possible that some of its assumptions can be derived from some of
the other assumptions.
This isn't desirable.
When applying DDC in practice it takes effort to gain strong
confidence that each assumption is true.
It would be better if we could either (1) prove that some assumptions are
derived from others
(in which case it no longer needs to be separate assumption),
or (2) prove that all the assumptions are independent of each other
in their respective proofs.

Given some set of consistent assumptions $Q$, and another assumption $P$, then $P$ is independent of $Q$ if both $Q \cup \{P\}$ and $Q \cup \{\neg P\}$ yield consistent assumptions. This is noted, for example, by Kunen~\cite[page xi]{kunen1980set}. It's already formally proved that our entire base set of DDC assumptions is consistent, so the first half of that definition ($Q \cup \{P\}$) is already proven. Therefore, to prove $P$'s independence, we only have to show that $Q \cup \{\neg P\}$ is consistent. In first-order logic, the easiest way to show that a set of assumptions is consistent is to find a model where they are all true. So, for each assumption $P$, we just need to find a model where $P$ is false but all other assumptions in $Q$ are true.

Hesseling explains that Heyting~\cite{heyting1930regeln} uses this exact approach to prove the independence of his formalization of intuitionist logic, and credits Bernays with the technique in which ``what we would nowadays call models in which each time ten of the eleven given axioms are true, the eleventh one false\dots\ [by] varying the eleventh axiom and presenting a different model to which the new combination of axioms applies, he proves that all the axioms are independent of each other''~\cite[page 278]{hesseling2003gnomes}.

To prove this formally, we used automated logic tools to establish the bounds of this independence. Specifically, to show that a set of $n$ assumptions is independent, we must prove that no single assumption can be derived from the remaining $n-1$ assumptions. In first-order logic, proving independence is not a matter of generating proofs of correctness; it is a search for counterexamples.

In our case, for each assumption $a_i$, we must show that there exists a mathematically consistent world where $a_i$ is false, yet all other assumptions in a given proof remain true. If we do this for all three proofs, we've shown that all of the assumptions are independent of each other.

We worked to accomplish this by feeding our formal axioms into Mace4 \cite{mccune2009prover9}, which searches for finite mathematical models. Mace4 is not at the cutting edge of model creation technology today. However, it's still a fine tool, and mace4 takes the same syntax as the prover9 tool which was used to generate the original proofs. Using the tool specifically designed to take the syntax avoided the risk of mistranslating the statements if we switched to a different system.

\subsection{The Mace4 Scoping Trap}

When attempting to formally prove independence using Mace4, early drafts ran into a subtle semantic roadblock while analyzing the second assumption of goal \path{always_equal}. Specifically, the issue occurred with the axiom defining a portable and deterministic program, labelled \path{define_portable_and_deterministic}.

Consider this attempted negation of the assumption, where we take the original aassertion and wrap it with \path{- ( ... )}:

\begin{verbatim}
formulas(assumptions).
  - ( ( portable_and_deterministic(Source,
          Language, Input) &
        exactly_correspond(Executable1, Source,
	  Language, Environment1) &
        exactly_correspond(Executable2, Source,
	  Language, Environment2)) ->
          ( converttext(run(Executable1, Input,
	      EnvEffects1, Environment1),
              Environment1, Target) =
            converttext(run(Executable2, Input,
	      EnvEffects2, Environment2),
              Environment2, Target)))
    # label(define_portable_and_deterministic).
end_of_list.
\end{verbatim}

At first glance, this looks completely reasonable. We want the rule to fail, so we put a negation symbol ($-$) in front of it, and we surround everything else with parenthesis to ensure it has the highest priority. Unfortunately, this creates a catastrophic structural error due to how Mace4 handles unquantified variables.

In Prover9 and Mace4, when you use capitalized terms (like \path{Source}, \path{Language}, or \path{Input}) without an explicit quantifier, the tool treats them as free variables. To process the file, the engine must bind these variables. It does this by automatically wrapping the entire expression in an implicit universal quantifier ($\forall$). Crucially, because the negation symbol ($-$) was written \emph{inside} the statement, Mace4 applied the implicit universal quantifier \emph{outside} the negation. Mathematically, Mace4 interpreted the assumption as:
$$\forall \vec{x} \, \neg (\text{Premise}(\vec{x}) \rightarrow \text{Conclusion}(\vec{x}))$$

For an implication ($\rightarrow$) to be false, the premise must be true \emph{and} the conclusion must be false. Because of the universal quantifier, we didn't ask Mace4 to find a single, specific loophole. Instead, we commanded it to construct a universe where \emph{every single object} in existence must satisfy the premise (meaning everything in existence is portable, deterministic, and exactly corresponding) and yet \emph{every single compiled output} must simultaneously fail to be equal to every other output. This defines a completely absurd, hyper-restrictive mathematical universe. The tool didn't fail because our DDC logic was flawed; it failed because we asked it to build an impossible world.

To prove independence, we don't want to say that the rule fails everywhere for everything. We want to prove that there is \emph{at least one} counterexample. Mathematically, a true counterexample requires an existential quantifier ($\exists$) wrapped around the negation: $$\exists \vec{x} \, (\text{Premise}(\vec{x}) \land \neg \text{Conclusion}(\vec{x}))$$

In raw Mace4 syntax, you would have to manually write out every single variable to force this behavior:
\begin{verbatim}
exists Source Language Input Executable1
       Environment1 Executable2 Environment2
  EnvEffects1 EnvEffects2 Target (
    portable_and_deterministic(Source, Language,
      Input) &
    exactly_correspond(Executable1, Source,
      Language, Environment1) &
    exactly_correspond(Executable2, Source,
      Language, Environment2) &
    - ( converttext(run(Executable1, Input,
          EnvEffects1, Environment1),
          Environment1, Target) =
        converttext(run(Executable2, Input,
	  EnvEffects2, Environment2),
          Environment2, Target))
).
\end{verbatim}

Writing this out manually is tedious and highly susceptible to typos or missed variables, which raised some concerns.

We then realized there was a much safer engineering-focused solution: let the tool do the negation for us, including handling the quantifiers. Instead of manually negating the rule and stuffing it into the assumptions list, we left the statement completely positive and placed it inside the \path{formulas(goals)} block:

\begin{verbatim}
formulas(assumptions).
    % All other DDC assumptions unmodified
end_of_list.

formulas(goals).
    % The assumption to be shown independent
end_of_list.
\end{verbatim}

When Prover9 or Mace4 sees a statement in the \path{goals} block, it utilizes a built-in automated mechanism that correctly negates the expression behind the scenes. By utilizing the \path{goals} block, we eliminate the risk of quantifier mismatch. Mace4 will search for a model that satisfies all our assumptions while breaking the goal in exactly \emph{one} valid place. If it finds that model, our independence proof is secure, our syntax remains clean, and we avoid subtle scoping traps.

\section{Success}

Once this was addressed, for each assumption, Mace4 successfully found a countermodel with a negated assumption. This formally guarantees that none of the DDC assumptions are redundant.

Since the assumptions are independent, application of DDC cannot take a lazy shortcut. To rely on the security guarantees of the DDC proof, you must establish the validity of every one of the proof's assumptions.

The specific models found can be viewed at \href{https://dwheeler.com/trusting-trust}{\path{https://dwheeler.com/trusting-trust}}.

\section{Conclusion and Next Steps}

By formally proving the independence of the assumptions underlying the three proofs of Diverse Double-Compiling (DDC), we have established that there are no redundant axioms in the DDC mathematical framework. When practitioners apply DDC to secure their binary distribution pipelines, they cannot assume that any single condition is a "freebie" implied by the others. Each assumption, from compiler determinism to the independence of the environment configurations, must be validated with care. On the other hand, they can be confident that their efforts are not wasted, since each assumption really does need to be verified.

The Mace4 scoping trap underscores that formal tools are only as robust as our understanding of their underlying semantics. Minor changes in the placement of negations can easily transform a standard logical inquiry into a mathematical impossibility.

This work lays a rigorous foundation, but it is just one step in securing the global digital supply chain. Ultimately, the goal of this research is to help counter attackers, even sophisticated attackers, so that all of us can use software that is increasingly trustworthy.

\bibliographystyle{ACM-Reference-Format}
\bibliography{\jobname}

\end{document}
